%==============================================================================
\section{Introduction --- Giới thiệu về học máy}
\label{sec:introduction}
%==============================================================================

\begin{dongco}[title={Bối cảnh: thời đại dữ liệu}]
Chúng ta đang sống trong “thời đại dữ liệu”: năng lực tính toán và lưu trữ tăng, dữ liệu tăng mỗi ngày.
Thách thức lớn là \textbf{làm cho dữ liệu có ý nghĩa} (make sense of data).
Doanh nghiệp/tổ chức xây hệ thống thông minh dựa trên Data Science, Data Mining và Machine Learning.
\end{dongco}

\begin{vidu}[title={Hình minh hoạ quy trình ML}]
\tsimg{00_khai_niem_nen/img04.png}
\end{vidu}

\subsection{Học máy hoạt động như thế nào?}

\begin{phuongphap}[title={3 thành phần chính}]
Cơ chế “máy học từ mô hình” được chia thành 3 phần:
\begin{itemize}
  \item \textbf{Decision Process}: dựa trên dữ liệu vào và nhãn đầu ra, mô hình hình thành một logic về pattern nhận diện được.
  \item \textbf{Cost Function}: đo lỗi giữa giá trị kỳ vọng và giá trị dự đoán; dùng để đánh giá hiệu năng.
  \item \textbf{Optimization Process}: tối thiểu hoá cost bằng cách điều chỉnh trọng số trong giai đoạn training; lặp đánh giá--tối ưu cho đến khi lỗi giảm.
\end{itemize}
\end{phuongphap}

\subsection{Vì sao cần học máy?}

\begin{dongco}[title={Quyết định dựa trên dữ liệu ở quy mô lớn}]
Con người rất thông minh, còn AI vẫn chưa vượt con người ở nhiều mặt.
Vậy tại sao vẫn cần “làm cho máy học”?
Một lý do phù hợp: \textbf{ra quyết định dựa trên dữ liệu} với \textbf{hiệu quả và quy mô}.
\end{dongco}

\begin{ynghia}[title={Tự động hoá và quyết định dựa trên dữ liệu}]
Gần đây, tổ chức đầu tư mạnh vào AI/ML/DL để khai thác thông tin then chốt từ dữ liệu cho các nhiệm vụ thực tế.
Các quyết định dựa trên dữ liệu có thể dùng thay cho việc lập trình logic ở các bài toán vốn khó lập trình bằng tay.
Ta vẫn cần trí tuệ con người, nhưng cũng cần giải quyết bài toán thực tế hiệu quả ở quy mô lớn.
\end{ynghia}

\subsection{Lược sử học máy}

\begin{phuongphap}[title={Các mốc}]
\begin{itemize}
  \item 1959: Arthur Samuel tạo chương trình ước tính xác suất thắng trong trò checkers.
  \item 1960--1970: giai đoạn neural networks phát triển.
  \item Sau đó ML tiến triển nhờ phương pháp thống kê như Bayesian networks, decision tree learning.
  \item 2010s: “cách mạng Deep Learning” với NLP, CNN, speech recognition.
  \item Hiện nay ML trở thành công nghệ phổ biến trong nhiều lĩnh vực (y tế, tài chính, giao thông, \ldots).
\end{itemize}
\end{phuongphap}

\subsection{Use cases của học máy}

\begin{phuongphap}[title={Một số use case}]
\begin{itemize}
  \item \textbf{Recommendation system}: engine gợi ý theo sở thích/engagement trước đó.
  \item \textbf{Voice assistants}: nhận dạng tiếng nói + xử lý ngôn ngữ + tổng hợp giọng.
  \item \textbf{Fraud detection}: phát hiện hoạt động bất thường (thường ở tài chính).
  \item \textbf{Health care}: chẩn đoán, cải thiện độ chính xác ảnh y khoa, cá nhân hoá điều trị.
  \item \textbf{RPA}: tự động hoá tác vụ lặp.
  \item \textbf{Drive-less cars}: xe tự lái; ví dụ Tesla.
  \item \textbf{Computer vision}: hiểu ảnh/video; nhận diện khuôn mặt.
\end{itemize}
\end{phuongphap}

\subsection{Ưu điểm và nhược điểm của học máy}

\subsubsection{Ưu điểm}
\begin{phuongphap}[title={Advantages of Machine Learning}]
\begin{itemize}
  \item \textbf{Tự động hoá}: tác vụ lặp có thể tự động hoá; ví dụ chatbot giảm thời gian chờ.
  \item \textbf{Cải thiện trải nghiệm và ra quyết định}: phân tích dữ liệu lớn để ra quyết định; cá nhân hoá sản phẩm/dịch vụ.
  \item \textbf{Ứng dụng rộng}: áp dụng trong nhiều lĩnh vực.
  \item \textbf{Cải tiến liên tục}: mô hình có thể học tiếp khi được retrain bằng dữ liệu mới.
\end{itemize}
\end{phuongphap}

\subsubsection{Nhược điểm}
\begin{phuongphap}[title={Disadvantages of Machine Learning}]
\begin{itemize}
  \item \textbf{Thu thập dữ liệu khó}: cần dữ liệu liên quan, không thiên lệch, chất lượng tốt.
  \item \textbf{Kết quả có thể không chính xác}: độ tin cậy phụ thuộc dữ liệu và thuật toán.
  \item \textbf{Dễ sai nếu dữ liệu/thuật toán sai}: ví dụ dữ liệu nhỏ làm mô hình học pattern kém, dự đoán lệch.
  \item \textbf{Cần bảo trì}: mô hình phải được giám sát và cập nhật liên tục.
\end{itemize}
\end{phuongphap}

\subsection{Thách thức và vấn đề đạo đức}

\begin{phuongphap}[title={Challenges in Machine Learning}]
\begin{itemize}
  \item \textbf{Data privacy}: dữ liệu có thể nhạy cảm; cần ẩn danh, bảo vệ dữ liệu, cân bằng sử dụng và quyền riêng tư.
  \item \textbf{Impact on jobs}: tự động hoá thay đổi việc làm; đồng thời tạo nghề mới (data scientist, ML engineer, \ldots).
  \item \textbf{Bias \& discrimination}: tránh dùng thuộc tính nhạy cảm sai cách (race, gender, \ldots).
  \item \textbf{Ethical consideration}: minh bạch, trách nhiệm giải trình, tác động xã hội.
\end{itemize}
\end{phuongphap}

\subsection{ML algorithms vs Traditional programming}

\begin{vidu}[title={So sánh theo tiêu chí}]
\begin{tabularx}{\linewidth}{>{\raggedright\arraybackslash}p{3.2cm}XX}
\textbf{Tiêu chí} & \textbf{Thuật toán học máy} & \textbf{Lập trình truyền thống} \\
\midrule
Cách tiếp cận & Máy học bằng cách huấn luyện trên dữ liệu lớn. & Lập trình viên viết quy tắc tường minh (code) để máy làm theo. \\
Dữ liệu & Phụ thuộc mạnh vào dữ liệu; dữ liệu quyết định chất lượng mô hình. & Phụ thuộc ít hơn; đầu ra chủ yếu phụ thuộc logic được mã hoá. \\
Độ phức tạp bài toán & Hợp bài toán phức tạp (image segmentation, NLP) cần nhận diện pattern. & Hợp bài toán có đầu ra/logic rõ ràng. \\
Tính linh hoạt & Linh hoạt, thích nghi khi retrain bằng dữ liệu mới. & Ít linh hoạt; thay đổi phải sửa code thủ công. \\
Kết quả & Khó dự đoán hoàn toàn vì phụ thuộc dữ liệu, mô hình, \ldots & Có thể dự đoán chính xác nếu logic đúng và đầy đủ. \\
\end{tabularx}
\end{vidu}

\subsection{ML vs Deep Learning; ML vs Generative AI}

\begin{ynghia}[title={ML vs Deep Learning}]
Deep learning là một nhánh của ML, khác nhau ở cách thuật toán học: deep learning dùng cấu trúc “giống não người” (mạng sâu).
Deep learning thường hiệu quả hơn ở bài toán phức tạp (ví dụ xe tự lái).
\end{ynghia}

\begin{ynghia}[title={ML vs Generative AI}]
Machine Learning chủ yếu dùng cho dự báo và ra quyết định; Generative AI tập trung tạo nội dung mới (ảnh/video/văn bản) theo pattern đã học.
\end{ynghia}

\subsection{Tương lai của học máy}

\begin{phuongphap}[title={Một số xu hướng}]
Nguồn đề cập các xu hướng như:
\begin{itemize}
  \item AutoML và synthetic data giúp ML dễ tiếp cận hơn.
  \item Quantum computing được kỳ vọng tăng tốc xử lý dữ liệu.
  \item Multi-modal AI: phân tích nhiều dạng input (text/speech/image/sensor).
  \item Edge computing, robotics, \ldots
\end{itemize}
\end{phuongphap}

\subsection{Cách học Machine Learning (5 bước)}

\begin{phuongphap}[title={Lộ trình 5 bước}]
\begin{enumerate}
  \item Học nền tảng (dữ liệu, thống kê, thuật toán, Python).
  \item Chọn framework (TensorFlow, PyTorch, scikit-learn, \ldots).
  \item Luyện với dữ liệu thật (Kaggle, UCI, \ldots).
  \item Tự làm dự án (từ đơn giản đến phức tạp).
  \item Tham gia cộng đồng (forum/meetup) để học hỏi và nhận phản hồi.
\end{enumerate}
\end{phuongphap}

